% Source fingerprint: d5311f1d77379e8582918b971818f50970a556f7480201a3a8a9e519ce628a85
% T07: raw TeX cells preserved; no numeric highlighting.
\begin{table}[htbp]
\centering\small\renewcommand{\arraystretch}{1.18}\setlength{\tabcolsep}{4pt}
\caption[Radon 与正则性的约定]{Radon 与正则性的约定。局部紧空间上的正则性须按正文条件证明。Fremlin 的完备、局部确定版本与本书 Borel 版本的差别仍以本节说明为准。}
\label{b1:tab:visual-t07}
\begin{tabularx}{\linewidth}{@{}>{\raggedright\arraybackslash}p{3.1cm}>{\raggedright\arraybackslash}p{4.7cm}>{\raggedright\arraybackslash}X@{}}
\toprule
条件或名称 & 本书的具体要求 & 未自动包含的性质 \\
\midrule
局部有限 & 每个点有有限测度的邻域 & 总质量有限、全空间可数穷竭 \\
紧内正则 & 每个 Borel 集的测度是其紧子集测度的上确界 & 开外正则、测度完备 \\
开外正则 & 每个 Borel 集的测度是其开超集测度的下确界 & 紧内正则 \\
Radon 测度 & 局部有限的 Borel 测度，且对全部 Borel 集紧内正则 & 不在定义中要求完备化或全局开外正则 \\
正则 Borel 测度 & 按本章约定同时具有紧内正则与开外正则 & 不同资料中同名用语可能要求较少或额外条件 \\
完备化 & 加入零测 Borel 集的任意子集，并延拓测度 & 它改变定义域；不能只因两测度在 Borel 集上一致便忽略此差别 \\
\bottomrule
\end{tabularx}
\end{table}
